\section{Introduction}\label{sec:introduction}

Chart-based models raise a basic question for return prediction. When a machine learns from a market chart, what has it learned? The answer shapes how we interpret predictability from recent price and volume data. Before the model sees a chart, the input window has already been selected, locally normalized, and rendered into a visual object. Strong chart CNN performance may therefore reflect the image layout, the local scaling of the input, or the ability of a flexible learner to extract nonlinear structure from price-volume information.

We study the representation question through three learned model families. The chart CNN learns from price and volume images, while CNN1D learns from scaled temporal sequences built from the same underlying data. The multimodal model combines the chart and sequence branches. CNN1D tests how much of the chart CNN signal can be recovered without the image layout, and the multimodal model tests whether the two representations contain complementary predictive information.

The sequence models are evaluated under three scaling schemes. Image scaling preserves the normalization embedded in chart construction, while cumulative-return and devolatized-return scaling provide alternative numerical representations of the same recent market history. We apply both a logistic classifier and CNN1D to these inputs. The logistic model shows how much predictive content is present in each encoding, and CNN1D shows what is added through nonlinear temporal learning.

The multimodal model provides a direct test of representation complementarity. If the chart and sequence branches extract distinct predictive information, combining them should improve the stock ranking beyond either branch alone. If their information mostly overlaps, multimodal gains should be limited. We complement this test with score correlations, encompassing regressions, and incremental explanatory-power tests because similar portfolio performance can arise from different stock rankings.

Our empirical analysis uses daily U.S. equity data from CRSP. Model development is confined to 1993--2000, and we report out-of-sample results for the full 2001--2024 period. We divide the evaluation period into a 2001--2019 baseline sample and a 2020--2024 extension sample. We compare the learned models with momentum (MOM), one-month short-term reversal (STR), one-week short-term reversal (WSTR), and a multi-horizon trend strategy (TREND). All methods use the same universe, timing discipline, portfolio construction rules, weighting schemes, and transaction-cost treatment.

The scaling choice is economically large. In the representative specification with a 20-day input and a five-day return horizon, the logistic classifier produces equal-weighted gross Sharpe ratios of 5.22 under image scaling, $-0.41$ under cumulative-return scaling, and 2.73 under devolatized-return scaling. Image scaling also produces the strongest logistic results across the 5-day and 60-day input windows. The image-scaled encoding therefore contains substantial predictive content before nonlinear temporal learning is added.

CNN1D models trained on image-scaled temporal sequences often match or outperform the chart CNN. At the five-day return horizon, CNN1D reaches equal-weighted gross Sharpe ratios of 6.24, 6.55, and 5.89 across the 5-day, 20-day, and 60-day input windows. The corresponding chart CNN values are 5.88, 6.52, and 3.72. The weekly signal therefore largely survives the move from chart image to temporal sequence under image scaling. At longer return horizons, performance is weaker and the additional contribution of nonlinear temporal learning is less consistent.

The multimodal evidence points to limited complementarity between the chart and sequence representations. Their scores have Spearman correlations of 0.73 and 0.67 in the two strongest weekly specifications, and their high-minus-low returns have correlations above 0.9 in the representative case. The multimodal model performs well, with a gross Sharpe ratio of 6.48 in the 20-day-input specification, compared with 6.52 for the chart CNN and 6.55 for CNN1D. Its limited improvement indicates that much of the predictive information is already present in the two single-branch models.

The longer-window weekly case shows stronger representation-specific variation. When the 60-day input is used to predict the five-day return, the chart CNN and CNN1D score correlation falls to 0.44. CNN1D also preserves its rankings more consistently across input windows than the chart CNN. The encompassing regressions and incremental explanatory-power tests identify CNN1D as the most consistent source of incremental predictive content, while the chart CNN retains independent content in selected specifications.

We compare the learned models with MOM, STR, WSTR, and TREND under common portfolio construction rules. In the full-sample equal-weighted weekly comparison, WSTR is the strongest gross benchmark with a Sharpe ratio of 2.86, followed by TREND at 2.49. The strongest learned specifications produce gross Sharpe ratios around 6.5. The learned-model advantage is clearest in equal-weighted weekly portfolios and narrows at longer horizons and under value-weighting.

Trading assumptions materially affect the economic interpretation. The representative full-universe equal-weighted weekly portfolios retain net Sharpe ratios near four after the size-dependent 10--20 basis-point transaction-cost adjustment. Value-weighted weekly portfolios are weak or negative after costs. The portfolios remain positive with much weaker performance after a price filter, while dollar-volume, no-microcap, and combined tradability filters produce negative net returns. The economically large signal is therefore concentrated in the smaller and less liquid part of the cross-section.

Performance weakens in the 2020--2024 extension sample. The representative learned models have equal-weighted net Sharpe ratios between 4.78 and 5.28 in the 2001--2019 baseline sample, compared with 1.11 to 1.35 in the extension sample. WSTR and TREND also weaken during 2020--2024. The decline therefore appears across learned models and price-based benchmarks, and the learned-model advantage narrows relative to the baseline sample. Several equal-weighted learned specifications still retain positive net performance.

\subsection{Related literature}\label{subsec:related_literature}

Our paper relates first to research on technical analysis and return predictability from price histories. \textcite{loFoundationsTechnicalAnalysis2000} formalize visual chart patterns through systematic statistical recognition. \textcite{reimagining_price_trends_jiang} replace prespecified patterns with a CNN trained directly on chart images. \textcite{murrayChartingMachines2024} use machine learning on numerical histories of monthly returns and find nonlinear predictive content distinct from standard momentum and reversal, including among the largest 500 stocks. The chart CNN replication of \textcite{kaczmarekErodingNonScalableCNN2025} places durability, microcap concentration, and implementation at the center of the image-based evidence. We study a different margin. Our tests hold recent price and volume histories fixed and ask which transformations of those histories explain the chart model's ranking.

The second connection is to machine learning in empirical asset pricing. Flexible methods can capture nonlinear interactions that linear models miss \parencite{guEmpiricalAssetPricing2020}, while recent evidence shows that carefully specified regressions can remain competitive in stock-return prediction \parencite{chengMachineLearningNecessity2025}. Our results connect these findings through input representation. Local scaling creates a strong linear ordering, and temporal convolution adds most consistently at the weekly horizon. Model complexity and data construction consequently enter the prediction problem jointly.

The third connection is to the economic evaluation of return predictors. Trading costs can reverse conclusions drawn from gross anomaly returns \parencite{novyMarxTaxonomyAnomalies2016}, and predictor performance commonly declines outside the discovery sample \parencite{mcleanAcademicResearchDestroy2016}. The concentration of weekly reversal among small and illiquid stocks provides a particularly close comparison for our short-horizon portfolios \parencite{chenMaxingOutReversals2025}. We combine these implementation tests with representation diagnostics, so the evidence distinguishes what the models learn from where the resulting portfolio can be traded.

Against this background, we make three contributions. First, we identify local high--low scaling as a concrete source of chart CNN performance and show that a continuous temporal model can reproduce the central weekly portfolio. Second, we connect representation choice to signal substitutability through rank correlations, encompassing regressions, and multimodal fusion. Third, we map the signal's economic scope using transparent price rules, transaction-cost estimates, weighting schemes, liquidity screens, factor regressions, and a post-2019 evaluation. Together, these tests identify the representation mechanism and place strict limits on its tradable scope.

The rest of the paper proceeds as follows. Section~\ref{sec:design} describes the data, representations, models, and portfolio design. Section~\ref{sec:mechanism_results} presents the scaling and temporal-learning results and studies overlap across learned signals. Section~\ref{sec:economic_limits} compares the models with transparent price rules and evaluates costs, tradability, and sample stability. Section~\ref{sec:conclusion} concludes. The online appendix reports architecture details, complete horizon grids, auxiliary fusion estimates, and supporting portfolio tables.
